% META: {"id":"1NSI-TC-EX-012","chapitre":"1NSI-TYPES-CONSTRUITS","type_objet":"exercice","capacites":["1NSI-TYPES-CONSTRUITS-C1"],"parcours":3,"competences":["programmer","modeliser"],"duree_min":15,"mode_creation":"ex_nihilo","sources_inspiration":[],"status":"approved","origin":"needs_review","status_history":["needs_review","audited","accepted","approved"],"release_acceptance":"RELEASE_OWNER_BATCH_ACCEPTANCE","acceptance_closure_digest":"sha256:9b3ccf9a81c5520fbb7e03b7b2d3e7bf2057a02834b6d908bf3be5f49f3b553f","acceptance_receipt_digest":"sha256:4fad86f0282e0fa4e0419cca1ad6379ae7af5a280c04a4a90880f90a1c044242"}
% BEGIN-VERIFY
% def enveloppe_convexe_1d(points):
%     if len(points) == 0:
%         return None
%     x_min = points[0][0]
%     x_max = points[0][0]
%     y_min = points[0][1]
%     y_max = points[0][1]
%     for x, y in points:
%         if x < x_min:
%             x_min = x
%         if x > x_max:
%             x_max = x
%         if y < y_min:
%             y_min = y
%         if y > y_max:
%             y_max = y
%     return ((x_min, y_min), (x_max, y_max))
% assert enveloppe_convexe_1d([(1, 2), (3, 0), (-1, 5)]) == ((-1, 0), (3, 5))
% assert enveloppe_convexe_1d([(0, 0)]) == ((0, 0), (0, 0))
% END-VERIFY
\begin{exercice}{1NSI-TC-EX-012}{3}{15}
On dispose d'une liste de tuples \lstinline{(x, y)} représentant les positions de
stations météo sur une carte.

Écrire une fonction \lstinline{enveloppe_convexe_1d(points)} qui renvoie un tuple de
deux tuples : le coin inférieur gauche \lstinline{(x_min, y_min)} et le coin
supérieur droit \lstinline{(x_max, y_max)} du rectangle englobant toutes les stations.

\textbf{Exemple :}
\begin{console}
>>> enveloppe_convexe_1d([(1, 2), (3, 0), (-1, 5)])
((-1, 0), (3, 5))
\end{console}

\textit{Indication :} parcourir tous les points pour trouver les extremums.
\end{exercice}
